\section{Foreign Desk: Swift \& Kotlin Clients}
\label{sec:swift-kotlin}

\subsection{Swift}
The Swift compiler generates plain \texttt{Foundation}-based code (falling
back to \texttt{FoundationNetworking} on Linux) --- no third-party
dependency pulled in. It produces type-safe request/response objects and
an action wrapper per action, mirroring every other Emi compiler:
\begin{lstlisting}
emi swift --path module.yaml --output ./Sources/Generated
\end{lstlisting}
\texttt{emi swift:dto} and \texttt{emi swift:headers} generate just the
DTOs or just the headers, when the full compiler's output isn't needed.
Each action carries its own metadata struct (name, URL, method) and a
strongly typed response, so calling code never hardcodes a route or
guesses a response shape:
\begin{lstlisting}
struct GetSinglePostActionMeta {
    let name: String = "GetSinglePostAction"
    let url: String = "/posts/:id"
    let method: String = "Get"
}
struct GetSinglePostActionResponse {
    let statusCode: Int
    let headers: [String: String]
    let payload: Data?
}
\end{lstlisting}

\subsection{Kotlin}
The Kotlin compiler targets Android and JVM projects from the same
definitions, using \href{https://github.com/Kotlin/kotlinx.serialization}{kotlinx.serialization}
for JSON and \href{https://square.github.io/okhttp/}{OkHttp} for
networking --- both need to be on the consuming project's classpath, and
nothing else does:
\begin{lstlisting}
emi kotlin --path module.yaml \
  --output ./src/main/kotlin/generated
\end{lstlisting}
One invocation compiles everything a module declares --- \texttt{dtos:},
\texttt{entities:} (each becomes a dto plus a set of CRUD actions), and
\texttt{actions:} --- into a flat directory of \texttt{.kt} files plus
shared runtime files. \texttt{emi kotlin:dto} and \texttt{emi
kotlin:headers} narrow the output the same way their Swift counterparts
do. \texttt{--pkg} sets the Kotlin \texttt{package} every generated file
is written under (default \texttt{unknownpackage}) --- it only matters
once more than one module compiles into the same project, so cross-module
dto/entity references resolve to the right package.

\begin{sidebar}{Shared shape across both compilers}
\begin{itemize}[leftmargin=1.1em,itemsep=1pt,topsep=2pt]
  \item \texttt{--path} / \texttt{--output} --- same convention as every
        Emi sub-compiler
  \item Without \texttt{--output}, generated virtual files print as JSON
        to the terminal for inspection before anything touches disk
  \item \texttt{--tags} --- comma-separated compile features, same
        mechanism as the Go and JS compilers
\end{itemize}
\end{sidebar}
